\clearpage\begingroup\SourceLeaf{409b}\ArithmeticHead{78}{L}{STÄCKEL, GAUSS ALS GEOMETER.}
\renewcommand{\thefootnote}{\arabic{footnote})}\setcounter{footnote}{1}
\begin{center}
\textbf{26.}\par\smallskip
\textls[100]{Kegelschnitte und Flächen zweiter Ordnung.}
\end{center}
\noindent\hspace*{\parindent}Auf die Lehre von den Kegelschnitten ist \textsc{Gauss} als Astronom immer
wieder geführt worden. Besonders eifrig hat er sich damit im Frühjahr 1843
beschäftigt. In dem Briefe an \textsc{Schumacher} vom 12. Mai 1843 erzählt er, dass
er »anfangs durch zufällige Umstände« seit vier bis sechs Wochen in einige
mathematische Spekulationen hineingezogen worden sei, »wo ich immer wieder
durch neue Aussichten in andere Richtungen gelenkt wurde und vieles erreicht,
vieles verfehlt habe. . . . Jene Spekulationen betrafen grossenteils weniger
neue Sachen als Durchführung mir eigentümlicher Methoden: zuletzt u. a.
mehreres sich auf die Kegelschnitte Beziehendes. Mir ist dabei wiederholt
in Erinnerung gekommen, wie ich vor einem halben Jahrhundert, als ich
zuerst \textsc{Newtons} \textit{Principia} las\footnote{\textsc{Gauss} hat sein Exemplar der \textit{Principia} im Jahre 1794 erworben.}, mehreres unbefriedigend fand, namentlich seine
an sich herrlichen Sätze die Kegelschnitte betreffend. Aber ich las immer
mit dem Gefühl, dass ich durch das Erlernte nicht \textls[100]{Herr} der Sache wurde;
besonders quälte mich die \textls[100]{gerade Linie}, mit deren Hilfe ein Kegelschnitt
beschrieben werden kann\footnote{Es handelt sich um die Konstruktion eines Kegelschnitts mittels zweier um ihre Scheitelpunkte drehbarer Winkel, deren eines Schenkelpaar sich auf einer Geraden schneidet, während der Durchschnittspunkt des anderen Schenkelpaares den Kegelschnitt beschreibt, I. \textsc{Newton}, \textit{Philosophiae naturalis principia mathematica}, London 1687, Liber I, sectio 5, Lemma 21. Vgl. auch C. \textsc{Maclaurin}, \textit{Geometria organica, sive descriptio linearum curvarum universalis}, London 1720, erster Abschnitt.}. . . . \textls[100]{Herr} des Gegenstandes ist man doch erst
dann, wenn man alle andern, diese magische gerade Linie betreffenden Fragen
beantworten kann; namentlich will man wissen, welche Relationen diese gerade
\clearpage\SourceLeaf{410}\ArithmeticHead{79}{R}{ELEMENTARE UND ANALYTISCHE GEOMETRIE.}\setcounter{footnote}{0}
\noindent Linie zu den Elementen des Kegelschnittes habe, ob man diese Elemente selbst
mit Leichtigkeit aus der Lage jener geraden Linie und der [gegebenen Punkte
des Kegelschnittes] ableiten könne. Verschiedenes dieser Art kann ich jetzt
recht artig ausrichten, ich weiss aber nicht, ob ich selbst das Ganze durchführen
kann, da andere Geschäfte mich nötigen abzubrechen« (W. VIII, S. 295).
\par
Was \textsc{Gauss} damals über seine Spekulationen niedergeschrieben hat, ist
aus dem Nachlass W. VIII, S. 341—344 abgedruckt. Die ihm eigentümliche
Methode war wieder die Benutzung komplexer Grössen; mittels dieses Verfahrens
hatte er schon etwa seit 1831 begonnen, die Kegelschnitte zu behandeln
(W. VIII, S. 339—340).
\par
\textsc{Gauss} hat die abgebrochene Arbeit nicht wieder aufgenommen, nicht aus
Mangel an Zeit, sondern weil ihm zufällig ein Buch in die Hände fiel, worin,
wie er am 15. Mai 1843 an \textsc{Schumacher} schreibt, »die Quintessenz der Lehre
von den Kegelschnitten in nucem gebracht ist«. Es war der schon 1827 erschienene
\textit{Barycentrische Calcul} von \textsc{Möbius}, ein Buch, das er, als es ihm 1828
vom Verfasser zugegangen war, »ohne viele Erwartung davon zu haben, zunächst
auf die Seite gelegt und später völlig vergessen hatte«, das aber, wie er jetzt
»mit grossem Vergnügen« fand, »auf dem leichtesten Wege zur Auflösung aller
dahin gehörigen Aufgaben führt« (W. VIII, S. 297)\footnote{Vgl. auch den Brief an \textsc{Schumacher} vom 19. Mai 1843, Br. G.-\textsc{Sch.} IV, S. 151.}.
\par
Dass \textsc{Gauss} sich in das Buch von \textsc{Möbius} vertieft hat, bezeugen auch die
aus dem Nachlass abgedruckten Notizen über das Pentagramma mirificum
(Fragment [11], W. VIII, S. 109—111) und über den Resultantencalcul (W.
VIII, S. 298).
\par
Wenn \textsc{Gauss} im \textit{Barycentrischen Calcul} die Quintessenz der Lehre von
den Kegelschnitten erblickt hat, so wird man daraus schliessen dürfen, dass
ihm die Untersuchungen \textsc{Poncelets}, \textsc{Steiners} und \textsc{Plückers} fremd geblieben
waren. Um seine Stellung zur neueren Geometrie zu bezeichnen, genügt es
daher nicht zu sagen, er habe die analytischen Methoden bevorzugt, man muss
vielmehr hinzufügen, dass er kein inneres Verhältnis zu den Auffassungen
gewonnen hat, die der projektiven Geometrie eigentümlich sind.
\par
Ebenfalls in das Jahr 1843 sind Auszüge zu setzen, die sich \textsc{Gauss} aus
zwei in den Pariser Comptes rendus vom 24. April 1843 erschienenen Noten
\clearpage\SourceLeaf{411}\ArithmeticHead{80}{L}{STÄCKEL, GAUSS ALS GEOMETER.}\setcounter{footnote}{0}
\noindent \textsc{Cauchys} gemacht hat. Sie stehen teils auf einer Notiztafel, die \textsc{Gauss} im
April 1840 von \textsc{Schumacher} zum Geschenk erhalten hatte (Br. G.-\textsc{Sch.} III,
S. 369) und die sich gegenwärtig im Besitz seines Enkels, Herrn C. \textsc{Gauss} in
Hameln, befindet, teils in dem Handbuch 19 Be, S. 254—257; auch die Notiz
über die Kreisschnitte, ebenda S. 253, hängt damit zusammen. Diese Auszüge
verdienen um so mehr Beachtung, als die im Nachlass vorhandenen Notizen
über Abhandlungen, die \textsc{Gauss} gelesen hatte, lediglich aus der Göttinger
Studienzeit (1795—1798) stammen; aus der späteren Zeit sind uns jedenfalls
keine Aufzeichnungen dieser Art erhalten, und auch in den Handbüchern
fehlen sie, abgesehen von dem soeben erwähnten Ausnahmefall. \textsc{Gauss} muss
also in den Sätzen von \textsc{Cauchy} etwas Besonderes gefunden haben, vielleicht
Berührungspunkte mit eigenen Untersuchungen.
\par
In der ersten Note\footnote{A. L. \textsc{Cauchy}, \textit{Mémoire sur la synthèse algébrique}, Comptes rendus, t. 16, Paris 1843, S. 867, Oeuvres, 1. série, t. 7, Paris 1892, S. 382.} betrachtet \textsc{Cauchy} eine ganze rationale Funktion
der rechtwinkligen Koordinaten eines Punktes der Ebene und zeigt, dass man
ihren Werten eine einfache geometrische Bedeutung beilegen kann; hiervon
werden Anwendungen auf die Kegelschnitte gemacht.
\par
Die zweite Note\footnote{A. L. \textsc{Cauchy}, \textit{Notes annexées au Rapport sur le Mémoire de M. Amyot}, Comptes rendus, ebenda, S. 885, Oeuvres, ebenda, S. 377.} enthält eine analytische Lösung der Aufgabe von \textsc{Amyot},
eine Fläche zweiter Ordnung als geometrischen Ort der Punkte darzustellen,
bei denen das Produkt der Entfernungen von zwei festen Ebenen zu dem
Quadrat der Entfernung von einem festen Punkte in einem gegebenen Verhältnis
steht\footnote{Ein ausführlicher Bericht über die der Pariser Akadsmie eingereichte Abhandlung \textsc{Amyots}: \textit{Nouvelle méthode de génération et de discussion des surfaces du second ordre} von \textsc{Cauchy} steht Comptes rendus ebenda, S. 783, Oeuvres, ebenda, S. 325; die Abhandlung \textsc{Amyots} ist abgedruckt in \textsc{Liouvilles} Journal, 2. série, t. 8, 1843, S. 163.}. \textsc{Cauchy} zeigt, dass die Aufgabe, abgesehen von der Ausziehung
gewisser Quadratwurzeln, auf eine Gleichung dritten Grades führt,
die mit der bekannten Gleichung für die reziproken Quadrate der Hauptachsen
übereinstimmt, ein Ergebnis, das man bei Heranziehung der Kreisschnitte
leicht bestätigen wird.
\par
Zu der Gleichung dritten Grades bemerkt \textsc{Gauss}: »Dies Resultat ist ganz
identisch mit meinem eigenen, vor 24 Jahren publizierten, was auch auf einem
\clearpage\SourceLeaf{412a}\ArithmeticHead{81}{R}{ELEMENTARE UND ANALYTISCHE GEOMETRIE.}\setcounter{footnote}{0}
\noindent besonderen Blatte steht«, und fügt die Buchstabenvertauschung hinzu, die
\textsc{Cauchys} Gleichung in die seinige überführe. Die Abhandlung, auf die er
sich bezieht, ist die 1818 erschienene \textit{Determinatio attractionis, quam . . . exercet
planeta . . .}, und zwar handelt es sich um die Formel [3] (W. III, S. 341).
Das besondere Blatt ist die Seite 114 des Handbuchs 19 Be; es enthält die
kubische Gleichung für die reziproken Quadrate der Hauptachsen genau in
der von \textsc{Gauss} angegebenen Bezeichnung. Damit stimmt, dass eine Notiz auf
S. 103 des Handbuchs das Datum des 20. Februar 1817 trägt. In anderer
Bezeichnungsweise findet sich die kubische Gleichung auf S. 166 desselben
Handbuches; diese etwa aus dem Jahre 1831 stammende Notiz ist W. II,
S. 307 abgedruckt. In geschichtlicher Beziehung sei noch bemerkt, dass die
kubische Gleichung schon 1812 von \textsc{Hachette} und \textsc{Petit}\footnote{\textsc{Hachette} und \textsc{Petit}, \textit{De l’équation qui a pour racines les carrés des demiaxes principaux d’une surface du second ordre}, Correspondance sur l’école polytechnique, t. 2, 1812, S. 324, 327.} angegeben war
und dass \textsc{Cauchy} sie 1826 abgeleitet hatte\footnote{A. L. \textsc{Cauchy}, \textit{Leçons sur les applications du calcul infinitésimal à la géométrie}, t. I, Paris 1826, S. 240; Oeuvres, 2. série, t. 5, S. 250.}.
\clearpage\endgroup
